% Provenance (never printed): see sources/outline_sources.md (Year 2: radiative forcing
% quantified; feedbacks and climate sensitivity; detection and attribution; scenarios; the
% remaining carbon budget; components of sea-level rise; physical science only). Every statement
% from IPCC AR6 WG1 (SPM; chapters 3, 5, 7 and 9; Annex III data), ledger rows in
% sources/ledger/book2.md (ipcc:*) and sources/ledger/book4.md (ar6*:*). Curves and numbers from
% figdata/bachelor-2/anthropogenic.py and the Book 4 second-writer scratch calculator calc/ch27.py.
\chapter{मानवजनित जलवायु परिवर्तन}\label{ch:b2:anthropogenic-climate-change}

1970 के बाद से वैश्विक सतह कम से कम पिछले दो हज़ार वर्षों की किसी भी अन्य पचास-वर्षीय अवधि से अधिक
तेज़ी से गर्म हुई है (\emph{उच्च विश्वास})। स्कूल की पुस्तकों ने बताया कि क्या बदला है और
क्यों। यह अध्याय उस वृत्तांत को एक बजट में बदलता है: मानवीय चालक प्रति वर्ग मीटर कितने वाट जोड़ते हैं,
गर्म होता ग्रह उनमें से कितने अंतरिक्ष को लौटाता है, कार्बन डाइऑक्साइड का हर टन कितना तापन लाता है,
और बढ़ता समुद्र गर्म होते जल और पिघलती बर्फ़ के बीच कैसे बँटा है। हर संख्या अपनी अनिश्चितता के साथ
आती है, मूल्यांकनों की अंशांकित भाषा में।
% ledger: ipcc:fastest-2000yr, ipcc:likelihood

\begin{recall}
\textbf{स्कूल की पुस्तकों से}: विकिरणी प्रणोदन, कार्बन डाइऑक्साइड के दोगुनीकरण के लिए $\qty{3.93}{W.m^{-2}}$,
और उसका लघुगणकीय नियम; जलवायु प्रतिपुष्टि; \omterm{def:b2:anthropogenic-climate-change:ecs}{साम्य जलवायु संवेदनशीलता}, सर्वश्रेष्ठ
अनुमान \qty{3}{\celsius}, \emph{संभावित} 2.5 से \qty{4}{\celsius}; \omterm{def:b2:anthropogenic-climate-change:da}{आरोपण}, उत्सर्जन परिदृश्य
और प्रक्षेपण; अंशांकित भाषा: \emph{संभावित} 66 से \qty{100}{\percent}, \emph{बहुत
संभावित} 90 से \qty{100}{\percent}, और एक परास $[x \text{ से } y]$ बहुत संभावित परास होता है। मानवीय
प्रभाव ने 1850--1900 और 2010--2019 के बीच जलवायु को \qty{1.07}{\celsius} गर्म किया है।
\textbf{भौतिकी की पुस्तकों से}: कृष्णिका और स्टेफ़न--बोल्ट्ज़मान नियम; तापीय प्रसार।
\textbf{जीव-विज्ञान की पुस्तकों से}: कार्बन चक्र, वायुवाहित अंश, कार्बन अवशोषक और जीवमंडल का
कार्बन बजट। \textbf{\cref{ch:b2:deep-time-climate} से}: धीमा कार्बन चक्र, जो
कार्बन डाइऑक्साइड के किसी अतिरिक्त को केवल लाखों वर्षों में ही वापस लेता है।
\end{recall}
% ledger: ipcc:erf-2xco2, ipcc:log-co2, ipcc:ecs, ipcc:likelihood, ipcc:attribution

\section{प्रणोदन का परिमाणन}\label{sec:b2:anthropogenic-climate-change:forcing}

\begin{definition}[प्रभावी विकिरणी प्रणोदन]\label{def:b2:anthropogenic-climate-change:erf}
किसी विक्षोभ --- किसी ग्रीनहाउस गैस में, एरोसॉलों में, सूर्य में परिवर्तन --- का \emph{प्रभावी विकिरणी
प्रणोदन}\index{प्रभावी विकिरणी प्रणोदन} (ERF) वायुमंडल के शीर्ष पर विकिरण के शुद्ध अधोमुखी प्रवाह में
उसके द्वारा उत्पन्न परिवर्तन है, \unit{W.m^{-2}} में, जब वायुमंडल, मेघ और स्थल-सतह उसके अनुसार
समायोजित हो चुके हों, पर वैश्विक सतही तापमान के अनुक्रिया देने से पहले। गिने गए समायोजन वे हैं जो
स्वयं कारक से होते हैं, सतही तापमान से स्वतंत्र
रूप से।
\end{definition}
% ledger: ar6ch7:erf-def

परिभाषा प्रणोदन और अनुक्रिया के बीच रेखा इस आधार पर खींचती है कि क्या सतही तापमान पर निर्भर करता है,
गति के आधार पर नहीं। कार्बन डाइऑक्साइड का तात्क्षणिक प्रणोदन गैस जोड़े जाने के ठीक बाद प्रवाह का परिवर्तन
है; समतापमंडल, जो फिर कुछ महीनों में ठंडा होता है, उसे पहली बार बदलता है; क्षोभमंडल, उसकी जल-वाष्प
और उसके मेघ उसे फिर बदलते हैं। कार्बन डाइऑक्साइड के दोगुनीकरण के लिए मूल्यांकन समतापमंडल के
समायोजित हो जाने पर \qty{3.75}{W.m^{-2}} पाता है, और क्षोभमंडल में ऐसे समायोजन जो आंशिक रूप से
निरस्त होते हैं: उसके तापन से \qty{-0.60}{W.m^{-2}}, जल-वाष्प से +0.22, मेघों से +0.45 और स्थल-सतह से
+0.11, कुल मिलाकर \qty{+0.18}{W.m^{-2}}।
ERF $3.75 + 0.18 = \qty{3.93}{W.m^{-2}}$ है, \qty{6}{\percent} अधिक, पिछले मूल्यांकन के
\qty{3.71}{W.m^{-2}} की तुलना में (\cref{exo:b2:anthropogenic-climate-change:1})।
% ledger: ar6ch7:2xco2, ar6ch7:erf-def

1750 से हर कारक का प्रणोदन वर्ष-दर-वर्ष पुनर्निर्मित किया गया है। नीचे का चित्र उसे दिखाता है:
2019 में कार्बन डाइऑक्साइड \qty{2.16}{W.m^{-2}} पर, अन्य ग्रीनहाउस गैसें --- मीथेन,
नाइट्रस ऑक्साइड, हैलोजनित गैसें और ओज़ोन --- मिलकर \qty{1.69}{W.m^{-2}} पर, और
एरोसॉल \qty{-1.06}{W.m^{-2}} पर। कुल मानवीय प्रणोदन \qty{2.72}{W.m^{-2}} है, और उसका बहुत
संभावित परास, 1.96 से \qty{3.48}{W.m^{-2}}, मुख्यतः एरोसॉलों के कारण चौड़ा है। 1750--2014 पर उनका ERF
$-1.3$ [$-2.0$ से $-0.6$] \unit{W.m^{-2}} आँका गया है (\emph{मध्यम विश्वास}),
जिसका अधिकांश भाग मेघों पर उनके प्रभाव से आता है, जो एरोसॉलों के अधिक बूँदें देने पर अधिक चमकीले हो जाते हैं।
ज्वालामुखी उद्गार छोटे ऋणात्मक उछाल जोड़ते हैं (\cref{ch:b2:eruptive-styles}): अभिलेख का सबसे गहरा,
1816 में \qty{-4.8}{W.m^{-2}}, 1815 में तंबोरा के उद्गार के बाद आता है।
% ledger: data:b4-ar6-erf, ipcc:erf-total, ipcc:erf-agents, ar6ch7:aerosol, usgs:msh-comparisons

\begin{omfigure}
\centering
\begin{tikzpicture}
\begin{axis}[omrecord, width=0.92\linewidth, height=6.2cm, xmin=1750, xmax=2020, ymin=-5.2,
  ymax=3.8, xlabel={वर्ष}, ylabel={1750 से ERF (\unit{W.m^{-2}})},
  xtick={1750,1800,...,2000}, /pgf/number format/1000 sep={}, xtick pos=bottom, ytick pos=left,
  tick label style={font=\scriptsize}, label style={font=\small},
  legend style={at={(0.5,-0.24)}, anchor=north, font=\scriptsize, draw=black!30,
  legend columns=3, /tikz/every even column/.append style={column sep=6pt}},
  legend cell align=left]
% computed in figdata/bachelor-2/anthropogenic.py part erf; ledger: data:b4-ar6-erf
\addplot[name path=hi, draw=none, forget plot] table[x=year, y=hi] {figdata/out/bachelor-2-anthropogenic-erf.dat};
\addplot[name path=lo, draw=none, forget plot] table[x=year, y=lo] {figdata/out/bachelor-2-anthropogenic-erf.dat};
\addplot[black!15, forget plot] fill between[of=hi and lo];
\addplot[black!45, line width=0.4pt] table[x=year, y=natural] {figdata/out/bachelor-2-anthropogenic-erf.dat};
\addlegendentry{प्राकृतिक (ज्वालामुखी, सूर्य)}
\addplot[omThm, line width=1pt] table[x=year, y=co2] {figdata/out/bachelor-2-anthropogenic-erf.dat};
\addlegendentry{कार्बन डाइऑक्साइड}
\addplot[omThm!50!orange, line width=1pt, densely dashed] table[x=year, y=other_ghg] {figdata/out/bachelor-2-anthropogenic-erf.dat};
\addlegendentry{अन्य ग्रीनहाउस गैसें और ओज़ोन}
\addplot[omDef, line width=1pt] table[x=year, y=aerosol] {figdata/out/bachelor-2-anthropogenic-erf.dat};
\addlegendentry{एरोसॉल}
\addplot[black, line width=1.4pt] table[x=year, y=anthro] {figdata/out/bachelor-2-anthropogenic-erf.dat};
\addlegendentry{कुल मानवीय (धूसर: 5--95\,\%)}
\draw[black!40] (axis cs:1750,0) -- (axis cs:2020,0);
\node[anchor=west, font=\scriptsize, text=black!60, inner sep=1pt] at (axis cs:1817.5,-4.75) {$\leftarrow$ तंबोरा};
\end{axis}
\end{tikzpicture}
\par\smallskip{\small 1750 से 2019 तक \omterm{def:b2:anthropogenic-climate-change:erf}{प्रभावी विकिरणी प्रणोदन}, 1750 के सापेक्ष, चालक के अनुसार।
कुल मानवीय प्रणोदन 2019 में \qty{2.72}{W.m^{-2}} तक पहुँचता है, बहुत संभावित परास छायांकित है।
डेटा: IPCC Sixth Assessment Report, Working Group~I, Annex~III extended data (CC BY 4.0)।}
\end{omfigure}
% values: anthro 2.7204 (1.9569-3.4819), co2 2.1563, other_ghg 1.6850, aerosol -1.0583, volcanic -4.8162 in 1816 (anthropogenic.py numbers)

\section{ऊर्जा बजट}\label{sec:b2:anthropogenic-climate-change:energy}

धनात्मक प्रणोदन ग्रह को ऊर्जा प्राप्त कराता है। गर्म होने पर वह अंतरिक्ष में अधिक विकिरण करता है, और
लाभ घटता है; जलवायु तभी स्थिर होती है जब दोनों संतुलित हों।

\begin{definition}[पृथ्वी का ऊर्जा असंतुलन; जलवायु प्रतिपुष्टि प्राचल]\label{def:b2:anthropogenic-climate-change:eei}
\emph{पृथ्वी का ऊर्जा असंतुलन}\index{पृथ्वी का ऊर्जा असंतुलन} $N$ वायुमंडल के शीर्ष पर जलवायु तंत्र में
ऊर्जा का शुद्ध प्रवाह है, पूरे ग्लोब पर और एक वर्ष या अधिक पर औसत किया हुआ,
\unit{W.m^{-2}} में: अवशोषित सूर्य-प्रकाश घटा बहिर्गामी अवरक्त। यह वह दर है जिससे महासागर,
स्थल, बर्फ़ और वायु ऊष्मा प्राप्त करते हैं, पृथ्वी की सतह के प्रति वर्ग मीटर।
\end{definition}

\begin{definition}[जलवायु प्रतिपुष्टि प्राचल]\label{def:b2:anthropogenic-climate-change:lambda}
\emph{जलवायु प्रतिपुष्टि प्राचल}\index{जलवायु प्रतिपुष्टि प्राचल} $\lambda$ वायुमंडल के शीर्ष पर शुद्ध
अधोमुखी प्रवाह में वैश्विक सतही तापन के प्रति अंश परिवर्तन है,
\unit{W.m^{-2}.\celsius^{-1}} में, जो उस तापन के प्रति जलवायु तंत्र की अनुक्रिया से होता है।
स्थिर जलवायु में $\lambda < 0$ होता है: तापन अंतरिक्ष को ऊर्जा की हानि बढ़ाता है।
\end{definition}

स्कूल की पुस्तक ने संतुलन को $\lambda$ के परिमाण, एक धनात्मक संख्या, के साथ लिखा था; यहाँ से आगे
चिह्न रखा जाता है, जैसा मूल्यांकनों में, जो इस प्राचल को $\alpha$ लिखते हैं।
% ledger: ar6ch7:energy-budget

\begin{proposition}[वैश्विक ऊर्जा संतुलन]\label{prop:b2:anthropogenic-climate-change:balance}
मान लीजिए $\Delta F$ \omterm{def:b2:anthropogenic-climate-change:erf}{प्रभावी विकिरणी प्रणोदन} है और $\Delta T$ किसी साम्य अवस्था के सापेक्ष वैश्विक सतही
तापन। यदि बहिर्गामी प्रवाह की अनुक्रिया $\Delta T$ में रैखिक है, तो
\[
  \Delta N = \Delta F + \lambda\,\Delta T .
\]
नए साम्य पर $\Delta N = 0$, अतः $\Delta T_{\text{साम्य}} = -\Delta F/\lambda$।
\end{proposition}
\begin{proof}
परिभाषा से प्रणोदन स्थिर सतही तापमान पर शुद्ध अधोमुखी प्रवाह को $\Delta F$ से बदलता है, और तापन
के प्रति अनुक्रिया उसे प्रथम कोटि तक $\lambda\,\Delta T$ से बदलती है;
असंतुलन उनका योग है। उसे शून्य रखने पर साम्य मिलता है।
\end{proof}

$N$ को मापने का अर्थ है तंत्र द्वारा प्राप्त ऊष्मा को मापना, और उसका लगभग सारा भाग महासागर रखता है।
1971 से 2018 तक जलवायु तंत्र ने \qty{435}{ZJ} (ज़ेटाजूल, \qty{e21}{J}) प्राप्त किए: महासागर ने
\qty{91}{\percent}, \qty{242}{ZJ} ऊपरी \qty{700}{m} में और उसके नीचे \qty{154}{ZJ}; स्थल ने
\qty{5}{\percent}, पिघलती बर्फ़ ने \qty{3}{\percent}, वायुमंडल ने \qty{1}{\percent}। माध्य
असंतुलन 1971--2018 पर \qty{0.57}{W.m^{-2}} था, और वह बढ़ा है: 2006--2018 पर \qty{0.79}{W.m^{-2}}
[0.52 से 1.06] (\cref{exo:b2:anthropogenic-climate-change:2})।
% ledger: ar6ch7:inventory, ipcc:heating

\begin{omfigure}
\centering
\begin{tikzpicture}[font=\small, >=Stealth]
\draw[fill=omDef!8, draw=omDef!60] (-3.2,0) rectangle (3.2,1.6);
\node[align=center] at (0,1.05) {जलवायु तंत्र};
\node[align=center, font=\scriptsize] at (0,0.5) {महासागर \qty{91}{\percent}, स्थल \qty{5}{\percent}, बर्फ़ \qty{3}{\percent}, वायु \qty{1}{\percent}};
\draw[->, line width=2.2pt, omThm] (-2.0,3.3) -- node[pos=0.12, left=3pt, align=right, font=\scriptsize] {प्रणोदन $\Delta F$} (-2.0,1.7);
\draw[->, line width=2.2pt, omDef] (2.0,1.7) -- node[pos=0.88, right=3pt, align=left, font=\scriptsize] {अनुक्रिया $\lambda\,\Delta T$} (2.0,3.3);
\draw[->, line width=1.4pt, black!70] (0,-0.1) -- node[right=3pt, align=left, font=\scriptsize] {संचित: $\Delta N = \Delta F + \lambda\,\Delta T$} (0,-1.0);
\draw[black!40, dashed] (-3.4,2.4) -- (3.4,2.4) node[right, font=\scriptsize, text=black!60] {वायुमंडल का शीर्ष};
\end{tikzpicture}
\par\smallskip{\small प्रणोदित जलवायु का ऊर्जा बजट (योजनामूलक)। प्रणोदन वायुमंडल के शीर्ष पर ऊर्जा जोड़ता
है, गर्म ग्रह की अनुक्रिया उसमें से कुछ हटाती है ($\lambda\,\Delta T$ ऋणात्मक
है), और अंतर संचित होता है, अधिकतर महासागर में। 1971 से 2018 तक प्राप्त ऊष्मा के
हिस्से।}
\end{omfigure}
% ledger: ar6ch7:inventory

\begin{method}[ऐतिहासिक ऊर्जा बजट से एक जलवायु संवेदनशीलता]\label{met:b2:anthropogenic-climate-change:energy-budget}
\begin{enumerate}
  \item दो अवधियाँ चुनिए, एक साम्य के निकट संदर्भ-अवधि और एक हाल की; प्रणोदन में परिवर्तन $\Delta F$,
        असंतुलन में $\Delta N$ और वैश्विक तापमान में $\Delta T$ लीजिए।
  \item प्रतिपुष्टि प्राचल के लिए संतुलन को हल कीजिए: $\lambda = (\Delta N - \Delta F)/\Delta T$।
  \item तब दोगुनीकरण के प्रति संवेदनशीलता $-F_{2\times}/\lambda$ है; प्रणोदन के समय का तापन,
        $F_{2\times}$ के अनुपात में मापित, $F_{2\times}\Delta T/\Delta F$ है।
\end{enumerate}
\end{method}

1850--1900 से 2006--2019 तक मूल्यांकन $\Delta N = \qty{0.59}{W.m^{-2}}$, $\Delta T =
\qty{1.03}{\celsius}$ और $\Delta F = \qty{2.20}{W.m^{-2}}$ लेता है। तब $\lambda = (0.59 - 2.20)/1.03 =
\qty{-1.56}{W.m^{-2}.\celsius^{-1}}$, और संवेदनशीलता $3.93/1.56 = \qty{2.5}{\celsius}$ है,
वही मान जो मूल्यांकन पाता है, 1.6 से \qty{4.8}{\celsius} के बहुत संभावित परास के साथ। पर यह
एक \emph{प्रभावी} संवेदनशीलता है, अब तक के तापन की प्रतिपुष्टियों के लिए। तापन वहाँ
सबसे प्रबल रहा है जहाँ प्रतिपुष्टियाँ सबसे अधिक ऋणात्मक हैं, पश्चिमी उष्णकटिबंधीय प्रशांत में, और
पूर्वी प्रशांत तथा दक्षिणी महासागर पर दुर्बल; जैसे-जैसे वह प्रतिरूप विकसित होता है, प्रतिपुष्टि प्राचल के
कम ऋणात्मक होने की अपेक्षा है (\emph{उच्च विश्वास}), अतः दीर्घकालिक संवेदनशीलता ऐतिहासिक
अनुमान से अधिक है (\cref{exo:b2:anthropogenic-climate-change:7})।
% ledger: ar6ch7:energy-budget, ar6ch7:pattern, ipcc:erf-2xco2
% values: lambda_energy -1.563, ecs_energy 2.514 (anthropogenic.py numbers)

\section{प्रतिपुष्टियाँ और संवेदनशीलता}\label{sec:b2:anthropogenic-climate-change:feedbacks}

\begin{definition}[प्लांक अनुक्रिया; जल-वाष्प, तापह्रास दर और मेघ प्रतिपुष्टियाँ]\label{def:b2:anthropogenic-climate-change:feedbacks}
प्रतिपुष्टि प्राचल को उन राशियों के योगदानों में बाँटा जाता है जो सतही तापमान के साथ बदलती हैं।
\emph{प्लांक अनुक्रिया}\index{प्लांक अनुक्रिया} वह अतिरिक्त अवरक्त है
जो अंतरिक्ष में उत्सर्जित होता है जब सतह और पूरा क्षोभमंडल समान मात्रा में गर्म होते हैं।
\emph{जल-वाष्प प्रतिपुष्टि}\index{जल-वाष्प प्रतिपुष्टि} तापमान के साथ वायुमंडलीय जल-वाष्प की वृद्धि
के कारण प्रवाह का परिवर्तन है। \emph{तापह्रास दर
प्रतिपुष्टि}\index{तापह्रास दर प्रतिपुष्टि} ऊँचाई के साथ असमान तापन के कारण प्रवाह का परिवर्तन है। \emph{मेघ
प्रतिपुष्टि}\index{मेघ प्रतिपुष्टि} मेघों की मात्रा, ऊँचाई और
चमक में परिवर्तन के कारण प्रवाह का परिवर्तन है।
\end{definition}

\begin{proposition}[प्लांक अनुक्रिया]\label{prop:b2:anthropogenic-climate-change:planck}
यदि पृथ्वी अपनी सतह के अवरक्त का एक नियत अंश $1 - g$ अंतरिक्ष में उत्सर्जित करती है, जहाँ वह अवरक्त
$\sigma T^4$ है और सतह का तापमान $T$, तो उसकी \omterm{def:b2:anthropogenic-climate-change:feedbacks}{प्लांक अनुक्रिया} है
\[
  \lambda_{\text P} = -(1 - g)\,4\sigma T^3 .
\]
$g = 0.4$ और $T = \qty{288}{K}$ के साथ, $\lambda_{\text P} = \qty{-3.25}{W.m^{-2}.K^{-1}}$।
\end{proposition}
\begin{proof}
बहिर्गामी प्रवाह $(1 - g)\sigma T^4$ है; $T$ के सापेक्ष उसका अवकलज $(1 - g)\,4\sigma
T^3$ है, और एक अतिरिक्त हानि शुद्ध अधोमुखी प्रवाह में एक ऋणात्मक योगदान है। संख्यात्मक रूप से $0.6
\times 4 \times \num{5.670e-8} \times 288^3 = 3.25$।
\end{proof}
% ledger: ar6ch7:planck, const:sigma, ipcc:budget-surface
% values: planck -3.2509 (anthropogenic.py numbers)

अंश $g$ प्रसामान्यीकृत ग्रीनहाउस प्रभाव है: वायुमंडल \qty{159}{W.m^{-2}} रखता है, उन
\qty{398}{W.m^{-2}} में से जो सतह उत्सर्जित करती है, $g \approx 0.4$। मूल्यांकन, विस्तृत
विकिरण-कोडों के साथ, $-3.22$ [$-3.4$ से $-3.0$] पाता है (\emph{उच्च विश्वास})। जल-वाष्प वैश्विक पैमाने पर
क्लॉसियस--क्लैपेरॉन नियम का अनुसरण करती है, और अधिक वाष्प अधिक अवरक्त अवशोषित करती है:
\omterm{def:b2:anthropogenic-climate-change:feedbacks}{जल-वाष्प प्रतिपुष्टि} प्रबल रूप से धनात्मक है, उपग्रहों से \qty{1.85}{W.m^{-2}.\celsius^{-1}}।
उष्णकटिबंध में ऊपरी क्षोभमंडल सतह से अधिक गर्म होता है, और फिर अधिक विकिरण करता है:
\omterm{def:b2:anthropogenic-climate-change:feedbacks}{तापह्रास दर प्रतिपुष्टि} वहाँ ऋणात्मक है, ध्रुवों के निकट धनात्मक, और वैश्विक स्तर पर लगभग $-0.50$। दोनों
का मूल्यांकन साथ-साथ किया जाता है, \num{1.30} [1.1 से 1.5], क्योंकि कोई मॉडल जो एक में डालता है, उसे वह
दूसरे से लेने की प्रवृत्ति रखता है। पिघलती हिम और बर्फ़ सतह को अधिक गहरे रंग का बना देती हैं, +0.35; मेघ \num{0.42}
[$-0.10$ से $0.94$] जोड़ते हैं, सबसे बड़ी अनिश्चितता। शुद्ध ऋणात्मक \omterm{def:b2:anthropogenic-climate-change:feedbacks}{मेघ प्रतिपुष्टि} \emph{बहुत असंभावित} है।
% ledger: ar6ch7:planck, ar6ch7:wv-lr, ar6ch7:table710, ipcc:cloud-negative, ipcc:budget-surface

\begin{omfigure}
\centering
\begin{tikzpicture}
\begin{axis}[xbar, width=0.78\linewidth, height=5.4cm, xmin=-3.7, xmax=1.8, ymin=0.4, ymax=6.6,
  ytick={1,...,6}, yticklabels={शुद्ध, जैव-भूभौतिक और रासायनिक, मेघ, सतही एल्बिडो,
  जल-वाष्प + तापह्रास दर, प्लांक}, y tick label style={font=\scriptsize},
  tick label style={font=\scriptsize}, xlabel={प्रतिपुष्टि प्राचल (\unit{W.m^{-2}.\celsius^{-1}})},
  label style={font=\small}, bar width=7pt, axis x line=bottom, axis y line=left,
  axis line style={-}, xmajorgrids, grid style={gray!22, line width=0.3pt}, extra x ticks={0},
  extra x tick style={grid style={black!60, line width=0.5pt}}]
% ledger: ar6ch7:table710 (central values and very likely ranges)
\addplot[fill=omDef!55, draw=omDef!80!black, error bars/.cd, x dir=both, x explicit,
  error bar style={black!70}] coordinates {
  (-3.22,6) += (0.22,0) -= (0.18,0)
  (-0.01,2) += (0.26,0) -= (0.26,0)
  (-1.16,1) += (0.65,0) -= (0.65,0)};
\addplot[fill=omThm!55, draw=omThm!80!black, error bars/.cd, x dir=both, x explicit,
  error bar style={black!70}] coordinates {
  (1.30,5) += (0.20,0) -= (0.20,0)
  (0.35,4) += (0.25,0) -= (0.25,0)
  (0.42,3) += (0.52,0) -= (0.52,0)};
\end{axis}
\end{tikzpicture}
\par\smallskip{\small IPCC द्वारा आँकी गई जलवायु प्रतिपुष्टियाँ (छठी मूल्यांकन रिपोर्ट, कार्यकारी
समूह~I, तालिका 7.10), केंद्रीय मान और बहुत संभावित परास। धनात्मक प्रतिपुष्टियाँ तापन को प्रवर्धित करती हैं
(लाल), ऋणात्मक उसे मंद करती हैं (नीला); शुद्ध उनका योग है।}
\end{omfigure}

\begin{definition}[साम्य जलवायु संवेदनशीलता]\label{def:b2:anthropogenic-climate-change:ecs}
\emph{साम्य जलवायु संवेदनशीलता}\index{साम्य जलवायु संवेदनशीलता} (ECS) नए साम्य पर वह वैश्विक
सतही तापन है जो वायु की कार्बन डाइऑक्साइड के उसके औद्योगिक-पूर्व स्तर से दोगुना होने से होता है,
हिमचादरों को स्थिर रखते हुए।
\end{definition}

\begin{proposition}[संवेदनशीलता और प्रतिपुष्टि]\label{prop:b2:anthropogenic-climate-change:ecs}
यदि प्रतिपुष्टि प्राचल स्थिर है, तो $\text{ECS} = -F_{2\times}/\lambda$। $F_{2\times} =
\qty{3.93}{W.m^{-2}}$ और ऊपर की प्रतिपुष्टियों के शुद्ध $\lambda = \qty{-1.16}{W.m^{-2}.\celsius^{-1}}$
के साथ, $\text{ECS} = \qty{3.4}{\celsius}$; अकेली \omterm{def:b2:anthropogenic-climate-change:feedbacks}{प्लांक अनुक्रिया} $3.93/3.22 =
\qty{1.2}{\celsius}$ देती।
\end{proposition}
\begin{proof}
साम्य पर $\Delta F =
F_{2\times}$ के साथ \cref{prop:b2:anthropogenic-climate-change:balance} लागू कीजिए।
\end{proof}
% ledger: ar6ch7:table710, ipcc:erf-2xco2
% values: ecs_process 3.388 (anthropogenic.py numbers)

प्रतिपुष्टियाँ साक्ष्य की एक धारा हैं; 1850 से तापन, ऊर्जा बजट और अतीत की जलवायु
अन्य धाराएँ हैं। मिलकर वे \qty{3}{\celsius} की आँकी गई ECS देती हैं, \emph{संभावित}
2.5 से \qty{4}{\celsius} (\emph{उच्च विश्वास}), \emph{बहुत संभावित} 2 से \qty{5}{\celsius}।
चूँकि ECS $1/\lambda$ के रूप में बदलती है, $\lambda$ पर सममित अनिश्चितता ECS का ऐसा परास देती है जो
ऊँचे मानों की ओर खिंचा होता है (\cref{exo:b2:anthropogenic-climate-change:12})।
% ledger: ipcc:ecs

\begin{definition}[क्षणिक जलवायु अनुक्रिया]\label{def:b2:anthropogenic-climate-change:tcr}
\emph{क्षणिक जलवायु अनुक्रिया}\index{क्षणिक जलवायु अनुक्रिया} (TCR) उस क्षण का वैश्विक सतही
तापन है जब वायु की कार्बन डाइऑक्साइड, अपने औद्योगिक-पूर्व स्तर से प्रति वर्ष \qty{1}{\percent} बढ़ते हुए,
दोगुनी हो जाती है, आरंभ के लगभग सत्तर वर्ष बाद।
\end{definition}

TCR, ECS से छोटी है क्योंकि दोगुनीकरण तक पहुँचने पर महासागर अभी भी ऊष्मा ले रहा होता है।
उसका सर्वश्रेष्ठ अनुमान \qty{1.8}{\celsius} है, \emph{संभावित} 1.4 से \qty{2.2}{\celsius}
(\emph{उच्च विश्वास}); अकेला ऐतिहासिक ऊर्जा बजट भी $3.93 \times 1.03/2.20 =
\qty{1.8}{\celsius}$ देता है।
% ledger: ar6ch7:tcr, ar6ch7:energy-budget
% values: tcr_energy 1.840 (anthropogenic.py numbers)

\section{संसूचन और आरोपण}\label{sec:b2:anthropogenic-climate-change:attribution}

\begin{definition}[संसूचन; आरोपण; जलवायु अंगुलिछाप]\label{def:b2:anthropogenic-climate-change:da}
किसी जलवायु परिवर्तन का \emph{संसूचन}\index{संसूचन (जलवायु परिवर्तन का)} यह दिखाना है कि
जलवायु एक परिभाषित सांख्यिकीय अर्थ में, उसकी आंतरिक परिवर्तनशीलता जो उत्पन्न कर सकती है उससे आगे, बदली है,
कारण बताए बिना। \emph{आरोपण}\index{आरोपण} किसी संसूचित परिवर्तन में कई कारणों के
योगदानों का मूल्यांकन है, एक निर्धारित संभाविता के साथ। \emph{जलवायु
अंगुलिछाप}\index{जलवायु अंगुलिछाप} किसी एक चालक के प्रति अनुक्रिया का दिक् और काल में प्रतिरूप है,
जैसा जलवायु मॉडल अनुकरण करते हैं, जिसके द्वारा प्रेक्षणों में उस चालक की पहचान होती है।
\end{definition}

भिन्न चालक भिन्न अंगुलिछापें छोड़ते हैं। अधिक ग्रीनहाउस गैस क्षोभमंडल को गर्म करती है और समतापमंडल को
ठंडा, क्योंकि अधिक मोटा अवरक्त कंबल ऊपरी परतों को अंतरिक्ष में अधिक विकिरण कराता है;
अधिक चमकीला सूर्य दोनों को गर्म करता। प्रेक्षण पहला प्रतिरूप दिखाते हैं: क्षोभमंडल कम से कम 1950 के दशक
से गर्म हुआ है, और यह \emph{लगभग निश्चित} है कि समतापमंडल ठंडा हुआ है। यह
\emph{बहुत संभावित} है कि मानवीय ग्रीनहाउस गैसें 1979 से क्षोभमंडलीय तापन की मुख्य चालक थीं,
और \emph{अत्यंत संभावित} है कि ओज़ोन परत का क्षय 1979 से 1990 के दशक के मध्य तक निचले समतापमंडल
के शीतलन का मुख्य चालक था।
% ledger: ar6ch3:upper-air

\begin{method}[इष्टतम अंगुलिछाप-विधि]\label{met:b2:anthropogenic-climate-change:fingerprinting}
\begin{enumerate}
  \item हर चालक के प्रति अनुक्रिया का अलग-अलग अनुकरण कीजिए: अंगुलिछापें $x_1, \dots, x_k$, दिक् और काल
        में मानों के सदिशों के रूप में।
  \item प्रेक्षणों को $y = \sum_i \beta_i x_i + \varepsilon$ के रूप में लिखिए, जहाँ $\varepsilon$
        आंतरिक परिवर्तनशीलता है, जिसके सांख्यिकीय गुण लंबे अप्रणोदित अनुकरणों से लिए जाते हैं।
  \item मापन-गुणकों $\beta_i$ का आकलन न्यूनतम वर्गों से कीजिए, परिवर्तनशीलता के व्युत्क्रम से भारित
        ताकि सबसे शोरयुक्त घटक सबसे कम गिने जाएँ, उनके विश्वास-अंतरालों के साथ।
  \item शून्य से सार्थक रूप से बड़े $\beta_i$ का अर्थ है कि चालक $i$ के प्रति अनुक्रिया
        संसूचित है; 1 को समाहित करने और 0 को बाहर रखने वाले अंतराल का अर्थ है कि अनुकरित अनुक्रिया
        प्रेक्षणों से संगत है, और परिवर्तन का आरोपण अंशतः उस चालक पर किया जा सकता है।
\end{enumerate}
\end{method}
% ledger: ar6ch3:fingerprint

सतही तापमान पर लागू करने पर, विधि आरंभ में स्मरण कराया गया \omterm{def:b2:anthropogenic-climate-change:da}{आरोपण} देती है: 1850--1900 और
2010--2019 के बीच प्रेक्षित \qty{1.06}{\celsius} में से मानवीय प्रभाव ने \emph{संभावित} रूप से
0.8 से \qty{1.3}{\celsius} उत्पन्न किया, सर्वश्रेष्ठ अनुमान \qty{1.07}{\celsius}; प्राकृतिक चालक और
आंतरिक परिवर्तनशीलता किसी भी दिशा में अंश के कुछ दसवें भाग से अधिक योगदान नहीं देते
(\cref{exo:b2:anthropogenic-climate-change:9,exo:b2:anthropogenic-climate-change:10})।
% ledger: ipcc:attribution, ipcc:observed-2010-2019

\section{परिदृश्य और संचयी उत्सर्जन}\label{sec:b2:anthropogenic-climate-change:scenarios}

\begin{definition}[साझा सामाजिक-आर्थिक मार्ग]\label{def:b2:anthropogenic-climate-change:ssp}
\emph{साझा सामाजिक-आर्थिक मार्ग}\index{साझा सामाजिक-आर्थिक मार्ग} (SSP) इस शताब्दी में जनसंख्या,
अर्थव्यवस्थाएँ और प्रौद्योगिकी कैसे विकसित हो सकती हैं, इसके पाँच आख्यानों में से एक है। मूल्यांकन के
उत्सर्जन परिदृश्य SSP$x$-$y$ नाम से चिह्नित हैं: $x$ मार्ग का नाम है, और $y$ वह
अनुमानित विकिरणी प्रणोदन, \unit{W.m^{-2}} में, जिस तक परिदृश्य 2100 में पहुँचता है।
\end{definition}
% ledger: ar6spm:ssp

पाँच मुख्य परिदृश्य SSP1-1.9 से SSP5-8.5 तक जाते हैं। भौतिक विज्ञान का मूल्यांकन उन्हें
दिया हुआ मानता है: वह न उनकी संभाविता तौलता है, न आख्यानों का मूल्यांकन करता है। हर एक के लिए जलवायु मॉडल
1850--1900 से ऊपर 2081--2100 के लिए एक प्रक्षेपित तापन देते हैं, \qty{1.4}{\celsius} (बहुत संभावित 1.0
से 1.8) से SSP1-1.9 के अंतर्गत, \qty{4.4}{\celsius} (3.3 से 5.7) तक SSP5-8.5 के अंतर्गत। प्रक्षेपण एक
सरल नियम साझा करते हैं।
% ledger: ar6spm:ssp, ipcc:proj-temp

\begin{definition}[संचयी CO$_2$ उत्सर्जन के प्रति क्षणिक जलवायु अनुक्रिया]\label{def:b2:anthropogenic-climate-change:tcre}
\emph{संचयी \ce{CO2} उत्सर्जन के प्रति क्षणिक जलवायु
अनुक्रिया}\index{संचयी CO$_2$ उत्सर्जन के प्रति क्षणिक जलवायु अनुक्रिया} (TCRE) औद्योगिक-पूर्व काल से
कुल उत्सर्जित कार्बन डाइऑक्साइड की प्रति इकाई वैश्विक सतही तापन है, \ce{CO2} के प्रति
\qty{1000}{Gt} पर \unit{\celsius} में।
\end{definition}

\begin{proposition}[तापन संचयी उत्सर्जन का अनुसरण करता है]\label{prop:b2:anthropogenic-climate-change:tcre}
अच्छे सन्निकटन तक कार्बन डाइऑक्साइड से होने वाला तापन कुल उत्सर्जित मात्रा के समानुपाती होता है,
उत्सर्जनों का मार्ग चाहे जो हो: $\Delta T \approx \text{TCRE} \times E$, जहाँ $E$ संचयी
उत्सर्जन है।
\end{proposition}
\admitted
% ledger: ar6spm:tcre, ar6ch5:tcre-linear

यह नियम स्पष्ट नहीं है। प्रणोदन सांद्रता के केवल लघुगणक के रूप में बढ़ता है, अतः हर
जोड़ा गया टन कम गर्म करता है; पर उत्सर्जनों के संचित होने के साथ महासागर ऊष्मा और कार्बन को कम-से-कम
दक्षता से लेता है, अतः हर टन का अधिक भाग वायु में रहता है और उसका तापन कम विलंबित होता है।
दोनों प्रभाव लगभग निरस्त होते हैं, और सभी अध्ययन सहमत हैं कि यही नियम का मूल है। संबंध
लगभग-रैखिक है (\emph{उच्च विश्वास}), और कार्बन डाइऑक्साइड का हर \qty{1000}{Gt} \emph{संभावित} रूप से
0.27 से \qty{0.63}{\celsius} तापन उत्पन्न करता है, सर्वश्रेष्ठ अनुमान \qty{0.45}{\celsius}, यानी
\qty{1.65}{\celsius} प्रति \qty{1000}{PgC} कार्बन (\cref{exo:b2:anthropogenic-climate-change:4})। 1850 से 2019 तक
$\num{2390} \pm \qty{240}{Gt}$ कार्बन डाइऑक्साइड उत्सर्जित हुई।
% ledger: ar6ch5:tcre-linear, ar6spm:tcre, ar6spm:budget

\begin{definition}[शेष कार्बन बजट]\label{def:b2:anthropogenic-climate-change:budget}
किसी तापमान-सीमा के लिए \emph{शेष कार्बन बजट}\index{शेष कार्बन बजट} वह संचयी कार्बन डाइऑक्साइड
है जो किसी निर्धारित तिथि से अभी भी उत्सर्जित की जा सकती है, जबकि वैश्विक तापन एक निर्धारित प्रायिकता के साथ
उस सीमा से नीचे रहे, अन्य मानवीय चालकों से होने वाले तापन को
ध्यान में रखते हुए।
\end{definition}

\begin{method}[TCRE से शेष कार्बन बजट]\label{met:b2:anthropogenic-climate-change:budget}
\begin{enumerate}
  \item सीमा में से पहले से पहुँचा मानवीय तापन घटाइए, जो 2010--2019 के लिए
        \qty{1.07}{\celsius} है।
  \item वह तापन घटाइए जो अन्य चालक तब तक जोड़ देंगे जब तक कार्बन डाइऑक्साइड उत्सर्जन
        शुद्ध शून्य तक पहुँचें।
  \item TCRE से भाग दीजिए: यह 2010--2019 के मध्य से बजट है।
  \item तब से हुए उत्सर्जन घटाइए, और उन धीमी पृथ्वी-तंत्र प्रतिपुष्टियों के लिए एक गुंजाइश जिन्हें
        मॉडल छोड़ देते हैं (पर्माफ़्रॉस्ट का पिघलना सबसे बड़ी है)।
\end{enumerate}
\end{method}

मूल्यांकन इस तरह पाँच घटकों को जोड़ता है: TCRE; ऐतिहासिक तापन; अन्य चालकों द्वारा लगभग 0.1 से
\qty{0.2}{\celsius} का अतिरिक्त तापन; वह तापन जो उत्सर्जन रुकने के बाद भी जारी रह सकता है, औसतन शून्य
आँका गया (\emph{संभावित} रूप से $\pm
\qty{0.19}{\celsius}$ के भीतर); और पृथ्वी-तंत्र प्रतिपुष्टियाँ, प्रति अंश लगभग \qty{26}{Gt} कार्बन डाइऑक्साइड।
\qty{1.5}{\celsius} के लिए, जो 2010--2019 से \qty{0.43}{\celsius} ऊपर है, 2020 के आरंभ से शेष बजट
\qty{500}{Gt} कार्बन डाइऑक्साइड है, \qty{50}{\percent} प्रायिकता पर, और
\qty{1350}{Gt} \qty{2}{\celsius} के लिए। अन्य चालकों के बारे में किए गए चुनाव इन मानों को
\qty{220}{Gt} या अधिक खिसकाते हैं; सप्ताहांत समस्या पहले मान को फिर से बनाती है।
% ledger: ar6ch5:budget-parts, ar6spm:budget

\begin{dated}{2025}
% ledger: data:co2-mlo, data:gistemp (values read from the datasets as fetched 2026-10-22; offset computed)
मौना लोआ पर वार्षिक माध्य कार्बन डाइऑक्साइड 2025 में \qty{427.35}{ppm} तक पहुँची, लघुगणकीय नियम से
$3.93\log_2(427.35/278) = \qty{2.44}{W.m^{-2}}$ का प्रणोदन। NASA के तापमान
अभिलेख में 2024 सबसे गर्म वर्ष है, 1880--1900 के माध्य से \qty{1.50}{\celsius} ऊपर, जो उस अभिलेख के
सबसे आरंभिक वर्ष हैं; 2025 \qty{1.40}{\celsius} पर रहा।
\end{dated}

\begin{omfigure}
\centering
\begin{tikzpicture}
\begin{axis}[omrecord, width=0.86\linewidth, height=5.8cm, xmin=0, xmax=4000, ymin=0, ymax=2.6,
  xlabel={1850 से संचयी \ce{CO2} उत्सर्जन (\unit{Gt})},
  ylabel={1850--1900 से तापन (\unit{\celsius})}, xtick={0,500,...,4000},
  /pgf/number format/1000 sep={\,}, xtick pos=bottom, ytick pos=left, clip=false,
  tick label style={font=\scriptsize}, label style={font=\small}]
% computed in figdata/bachelor-2/anthropogenic.py part tcre; ledger: ar6spm:tcre, ar6spm:budget
\addplot[name path=hi, draw=none, forget plot] table[x=E_GtCO2, y=dT_hi] {figdata/out/bachelor-2-anthropogenic-tcre.dat};
\addplot[name path=lo, draw=none, forget plot] table[x=E_GtCO2, y=dT_lo] {figdata/out/bachelor-2-anthropogenic-tcre.dat};
\addplot[omThm!18, forget plot] fill between[of=hi and lo];
\addplot[omThm, line width=1.3pt] table[x=E_GtCO2, y=dT_mid] {figdata/out/bachelor-2-anthropogenic-tcre.dat};
\draw[densely dashed, black!60] (axis cs:0,1.5) -- (axis cs:4000,1.5);
\draw[densely dashed, black!60] (axis cs:0,2.0) -- (axis cs:4000,2.0);
\node[anchor=north west, font=\scriptsize, text=black!70] at (axis cs:0,1.5) {\qty{1.5}{\celsius}};
\node[anchor=north west, font=\scriptsize, text=black!70] at (axis cs:0,2.0) {\qty{2}{\celsius}};
\fill[black] (axis cs:2390,1.07) circle (2.2pt);
\draw[black, line width=0.8pt] (axis cs:2150,1.07) -- (axis cs:2630,1.07);
\draw[black, line width=0.8pt] (axis cs:2390,0.8) -- (axis cs:2390,1.3);
\node[anchor=north west, font=\scriptsize, align=left] at (axis cs:2440,1.02) {1850--2019:\\\qty{2390}{Gt}, \qty{1.07}{\celsius}};
\draw[<->, omDef, thick] (axis cs:2390,1.62) -- (axis cs:2890,1.62);
\node[anchor=south, font=\scriptsize, text=omDef, align=center] at (axis cs:2640,1.64) {AR6: \qty{500}{Gt}};
\node[anchor=north west, font=\scriptsize, text=omThm, align=left] at (axis cs:150,2.5) {TCRE का संभावित परास,\\0.27 से \qty{0.63}{\celsius} प्रति \qty{1000}{Gt}};
\end{axis}
\end{tikzpicture}
\par\smallskip{\small आँकी गई TCRE के लिए संचयी कार्बन डाइऑक्साइड उत्सर्जन के सापेक्ष तापन,
\qty{0.45}{\celsius} प्रति \qty{1000}{Gt} (रेखा) और उसका संभावित परास (छायांकित), मानवीय तापन
और अब तक के उत्सर्जन के साथ (बिंदु, संभावित परास पट्टियों के रूप में)। तीर 2020 से
\qty{1.5}{\celsius} के लिए, \qty{50}{\percent} पर, शेष बजट चिह्नित करता है: अकेली रेखा पढ़ने से मिलने वाले मान से छोटा,
क्योंकि अन्य चालक भी गर्म करते हैं।}
\end{omfigure}


\section{समुद्र-तल}\label{sec:b2:anthropogenic-climate-change:sea-level}

\begin{definition}[समुद्र-तल बजट; तापस्थैतिक समुद्र-तल वृद्धि]\label{def:b2:anthropogenic-climate-change:sea-level}
\emph{समुद्र-तल बजट}\index{समुद्र-तल बजट} वैश्विक माध्य समुद्र-तल के प्रेक्षित परिवर्तन की उसके
अनुमानित योगदानों के योग से तुलना है: गर्म होते महासागर का प्रसार, हिमनदों और ग्रीनलैंड तथा
अंटार्कटिक हिमचादरों का पिघलना, और स्थल पर संचित जल का परिवर्तन। \emph{तापस्थैतिक
समुद्र-तल वृद्धि}\index{तापस्थैतिक समुद्र-तल वृद्धि} वह भाग है जो स्थिर द्रव्यमान पर समुद्री जल के
तापीय प्रसार के कारण होता है।
\end{definition}

\begin{proposition}[किसी परत का तापीय प्रसार]\label{prop:b2:anthropogenic-climate-change:thermosteric}
मोटाई $H$ और तापीय प्रसार-गुणांक $\alpha$ वाली महासागर की एक परत, जो स्थिर द्रव्यमान और क्षेत्रफल पर
$\Delta T$ से समान रूप से गर्म होती है, $\Delta h = \alpha\,\Delta T\,H$ ऊपर उठती है।
\end{proposition}
\begin{proof}
एक स्तंभ का द्रव्यमान $\rho H A$, जिसका क्षेत्रफल $A$ है, स्थिर है, अतः प्रथम कोटि तक $\Delta H/H = -\Delta\rho/\rho =
\alpha\,\Delta T$, जहाँ $\alpha = -(1/\rho)\,\partial\rho/\partial T$।
\end{proof}
% ledger: teos10:alpha-beta

वास्तविक महासागर में $\alpha$ तापमान और दाब के साथ बढ़ता है, और तापन सतह के निकट केंद्रित
होता है। मूल्यांकन परिणाम को \qty{0.113}{m} प्रति \qty{e24}{J} (महासागर में जोड़ी गई ऊष्मा) की प्रसार-दक्षता
के रूप में सारांशित करता है, यानी पूरे महासागर के प्रति अंश तापन पर \qty{0.617}{m}
(\emph{मध्यम विश्वास})। उसकी \qty{3682}{m} की माध्य गहराई पर, यह एक प्रभावी $\alpha$ है,
$0.617/3682 = \qty{1.68e-4}{\celsius^{-1}}$ (\cref{exo:b2:anthropogenic-climate-change:8})।
% ledger: ar6ch9:expansion, charette:ocean-volume

बजट, जो नीचे के चित्र में खींचा गया है, अपनी अनिश्चितताओं के भीतर बंद होता है। 1971--2018 पर योगदानों का योग
\qty{2.00}{mm.yr^{-1}} है, \qty{2.33}{mm.yr^{-1}} की प्रेक्षित वृद्धि के सामने; 2006--2018 पर उनका योग
3.61 है, 3.69 के सामने। तापीय प्रसार सबसे बड़ा अकेला पद था, 1971--2018 पर योग का आधा, पर बर्फ़
आगे बढ़ रही है: 2006--2018 पर हिमनदों और हिमचादरों ने मिलकर \qty{1.62}{mm.yr^{-1}} दिए, जो
प्रसार से अधिक है। सतह के स्थिर होने के बहुत बाद तक महासागर गहराई में गर्म होता रहता है, और हिमचादरें
सदियों में अनुक्रिया देती हैं; समुद्र-तल का सदियों से सहस्राब्दियों तक बढ़ना तय है और वह हज़ारों वर्षों तक
ऊँचा रहेगा (\emph{उच्च विश्वास})। अगले दो हज़ार वर्षों में उसके लगभग 2 से
\qty{3}{m} बढ़ने का प्रक्षेपण है यदि तापन \qty{1.5}{\celsius} पर रहे, 2 से
\qty{6}{m} \qty{2}{\celsius} पर, और 19 से \qty{22}{m} \qty{5}{\celsius} पर (\emph{निम्न विश्वास})।
% ledger: ar6ch9:table95, ar6spm:sea-commit

\begin{omfigure}
\centering
\begin{tikzpicture}
\begin{axis}[ybar stacked, width=0.72\linewidth, height=6.2cm, ymin=0, ymax=4.3, bar width=26pt,
  xtick=data, symbolic x coords={1971--2018, 1993--2018, 2006--2018}, enlarge x limits=0.25,
  ylabel={वृद्धि की दर (\unit{mm.yr^{-1}})}, tick label style={font=\scriptsize},
  label style={font=\small}, ytick pos=left, xtick pos=bottom,
  legend style={at={(1.03,0.5)}, anchor=west, font=\scriptsize, draw=black!30},
  legend cell align=left, reverse legend]
% ledger: ar6ch9:table95
\addplot[fill=omDef!70, draw=omDef!90!black] coordinates {(1971--2018,1.01) (1993--2018,1.31) (2006--2018,1.39)};
\addlegendentry{तापीय प्रसार}
\addplot[fill=omDef!30, draw=omDef!70!black] coordinates {(1971--2018,0.44) (1993--2018,0.55) (2006--2018,0.62)};
\addlegendentry{हिमनद}
\addplot[fill=omThm!55, draw=omThm!80!black] coordinates {(1971--2018,0.25) (1993--2018,0.43) (2006--2018,0.63)};
\addlegendentry{ग्रीनलैंड}
\addplot[fill=omThm!25, draw=omThm!70!black] coordinates {(1971--2018,0.14) (1993--2018,0.25) (2006--2018,0.37)};
\addlegendentry{अंटार्कटिका}
\addplot[fill=black!20, draw=black!50] coordinates {(1971--2018,0.15) (1993--2018,0.31) (2006--2018,0.60)};
\addlegendentry{स्थलीय जल}
\node[fill=black, rotate=45, minimum size=5pt, inner sep=0pt] at (axis cs:1971--2018,2.33) {};
\node[fill=black, rotate=45, minimum size=5pt, inner sep=0pt] at (axis cs:1993--2018,3.25) {};
\node[fill=black, rotate=45, minimum size=5pt, inner sep=0pt] at (axis cs:2006--2018,3.69) {};
\end{axis}
\end{tikzpicture}
\par\smallskip{\small \omterm{def:b2:anthropogenic-climate-change:sea-level}{समुद्र-तल बजट}: वैश्विक माध्य समुद्र-तल की वृद्धि में प्रेक्षित योगदान, 2018 में
समाप्त होने वाली तीन अवधियों पर दरों के रूप में, एक के ऊपर एक रखे हुए, और प्रेक्षित दर (हीरे)।
ढेर और हीरे के बीच का अंतर वह भाग है जिसे बजट बंद नहीं करता। मान: IPCC छठी
मूल्यांकन रिपोर्ट, कार्यकारी समूह~I, तालिका 9.5।}
\end{omfigure}

\section*{अभ्यास}\label{sec:b2:anthropogenic-climate-change:exercises}

अभ्यासों के आँकड़े, अन्यथा न कहा जाए तो: दोगुनी कार्बन डाइऑक्साइड का प्रणोदन $F_{2\times} =
\qty{3.93}{W.m^{-2}}$, $C_0 = \qty{278}{ppm}$ से; स्टेफ़न--बोल्ट्ज़मान नियतांक $\sigma =
\qty{5.670e-8}{W.m^{-2}.K^{-4}}$; पृथ्वी की त्रिज्या \qty{6371}{km}; एक वर्ष
\qty{3.156e7}{s}; कार्बन और कार्बन डाइऑक्साइड के मोलर द्रव्यमान 12.011 और \qty{44.009}{g.mol^{-1}};
अध्याय की प्रतिपुष्टियाँ और \omterm{def:b2:anthropogenic-climate-change:sea-level}{समुद्र-तल बजट}।
% ledger: ipcc:erf-2xco2, const:sigma, earth:Rmean, ciaaw:abridged-2024, ar6ch7:table710, ar6ch9:table95

\begin{exercise}[$\star$]\label{exo:b2:anthropogenic-climate-change:1}
दोगुनी कार्बन डाइऑक्साइड का प्रणोदन समतापमंडल के समायोजित हो जाने पर \qty{3.75}{W.m^{-2}} है;
फिर क्षोभमंडल $-0.60$ (तापमान), +0.22 (जल-वाष्प), +0.45 (मेघ) और
+0.11 (स्थल-सतह) \unit{W.m^{-2}} से समायोजित होता है। (क)~ERF ज्ञात कीजिए। (ख)~वह पिछले मूल्यांकन के
\qty{3.71}{W.m^{-2}} से कितने प्रतिशत अधिक है? (ग)~कौन-सा समायोजन ऋणात्मक है, और एक गर्म
क्षोभमंडल प्रणोदन को क्यों घटाता है?
\end{exercise}
% ledger: ar6ch7:2xco2

\begin{exercise}[$\star$]\label{exo:b2:anthropogenic-climate-change:2}
जलवायु तंत्र ने 1971 से 2018 तक \qty{434.9}{ZJ} और 2006 से 2018 तक \qty{152.4}{ZJ} प्राप्त किए,
पहले और अंतिम वर्षों के मध्य-बिंदुओं के बीच गिनकर। (क)~पृथ्वी की सतह का क्षेत्रफल ज्ञात कीजिए।
(ख)~हर लाभ को \unit{W.m^{-2}} में एक माध्य असंतुलन में बदलिए। (ग)~महासागर ने
\qty{91}{\percent} लिया: 1971--2018 पर कितने ज़ेटाजूल?
\end{exercise}
% ledger: ar6ch7:inventory, earth:Rmean

\begin{exercise}[$\star$]\label{exo:b2:anthropogenic-climate-change:3}
(क)~चित्र की प्रतिपुष्टियों के केंद्रीय मानों का योग कीजिए, और शुद्ध मान से तुलना कीजिए। (ख)~शुद्ध मान
के साथ ECS ज्ञात कीजिए, और अकेली \omterm{def:b2:anthropogenic-climate-change:feedbacks}{प्लांक अनुक्रिया} के साथ तापन। (ग)~अन्य प्रतिपुष्टियाँ प्लांक तापन को
किस गुणक से प्रवर्धित करती हैं?
\end{exercise}

\begin{exercise}[$\star$]\label{exo:b2:anthropogenic-climate-change:4}
(क)~\qty{0.45}{\celsius} प्रति \qty{1000}{Gt} कार्बन डाइऑक्साइड की TCRE को कार्बन के प्रति
\qty{1000}{PgC} अंशों में बदलिए, और उसके संभावित परास, 0.27 से 0.63, के लिए भी यही कीजिए। (ख)~1850 से 2019 तक
उत्सर्जित \qty{2390}{Gt} कार्बन डाइऑक्साइड को पेटाग्राम कार्बन में बदलिए। (ग)~इन उत्सर्जनों के लिए
TCRE कितना तापन देती है?
\end{exercise}
% ledger: ar6spm:tcre, ar6spm:budget

\begin{exercise}[$\star\star$]\label{exo:b2:anthropogenic-climate-change:5}
(क)~सतह \qty{398}{W.m^{-2}} उत्सर्जित करती है और ग्रीनहाउस प्रभाव उसमें से \qty{159}{W.m^{-2}}
रखता है। $g$ ज्ञात कीजिए, और \qty{398}{W.m^{-2}} उत्सर्जित करने वाली कृष्णिका का सतही तापमान। (ख)~\qty{288}{K}
पर \omterm{def:b2:anthropogenic-climate-change:feedbacks}{प्लांक अनुक्रिया} ज्ञात कीजिए, और \qty{298}{K} पर, उसी $g$ के साथ। (ग)~क्या अकेली \omterm{def:b2:anthropogenic-climate-change:feedbacks}{प्लांक
अनुक्रिया} एक गर्म पृथ्वी को अधिक संवेदनशील बनाएगी या कम?
\end{exercise}
% ledger: ipcc:budget-surface, ar6ch7:planck

\begin{exercise}[$\star\star$]\label{exo:b2:anthropogenic-climate-change:6}
उपग्रह $1.85 \pm 0.32$ की \omterm{def:b2:anthropogenic-climate-change:feedbacks}{जल-वाष्प प्रतिपुष्टि} देते हैं और अध्ययन $-0.50 \pm 0.20$ \unit{W.m^{-2}.\celsius^{-1}}
की \omterm{def:b2:anthropogenic-climate-change:feedbacks}{तापह्रास दर प्रतिपुष्टि} (एक मानक विचलन)। (क)~उनका योग ज्ञात कीजिए।
(ख)~यदि दोनों स्वतंत्र होतीं, तो योग का मानक विचलन क्या होता? (ग)~मॉडल
योग को केवल \num{0.10} का मानक विचलन देते हैं। यह दोनों त्रुटियों के बारे में क्या बताता है?
\end{exercise}
% ledger: ar6ch7:wv-lr

\begin{exercise}[$\star\star$]\label{exo:b2:anthropogenic-climate-change:7}
1850--1900 से 2006--2019 तक, $\Delta N = \qty{0.59}{W.m^{-2}}$, $\Delta T = \qty{1.03}{\celsius}$
और $\Delta F = \qty{2.20}{W.m^{-2}}$। (क)~$\lambda$ और प्रभावी संवेदनशीलता ज्ञात कीजिए। (ख)~तापन के
प्रतिरूप के विकसित होने के साथ प्रतिपुष्टि प्राचल के $\lambda' = \qty{0.5}{W.m^{-2}.\celsius^{-1}}$ कम
ऋणात्मक होने की अपेक्षा है। संवेदनशीलता $-F_{2\times}/(\lambda + \lambda')$ ज्ञात कीजिए।
(ग)~$\lambda' = 0.1$ के साथ (ख) फिर से कीजिए।
\end{exercise}
% ledger: ar6ch7:energy-budget, ar6ch7:pattern, ar6ch7:pattern-ecs

\begin{exercise}[$\star\star$]\label{exo:b2:anthropogenic-climate-change:8}
(क)~पूरे महासागर के प्रति अंश तापन पर \qty{0.617}{m} की वृद्धि से, माध्य गहराई
\qty{3682}{m}, प्रभावी प्रसार-गुणांक ज्ञात कीजिए। (ख)~महासागर ने 1971--2018 पर \qty{396}{ZJ}
प्राप्त किए। \qty{0.113}{m} प्रति \qty{e24}{J} की प्रसार-दक्षता के साथ, तापस्थैतिक
वृद्धि ज्ञात कीजिए, और बजट के \qty{47.5}{mm} से तुलना कीजिए। (ग)~\qty{700}{m} की एक परत
\qty{0.2}{\celsius} गर्म होती है, $\alpha = \qty{2e-4}{\celsius^{-1}}$ के साथ: वह समुद्र-तल को कितना ऊपर उठाती है?
\end{exercise}
% ledger: ar6ch9:expansion, charette:ocean-volume, ar6ch7:inventory, ar6ch9:table95

\begin{exercise}[$\star\star$]\label{exo:b2:anthropogenic-climate-change:9}
चार दशकों में एक मॉडल अंगुलिछाप $x = (0.10, 0.20, 0.30, 0.40)$ \unit{\celsius} है; प्रेक्षण
$y = (0.12, 0.18, 0.33, 0.41)$ \unit{\celsius} हैं। (क)~न्यूनतम-वर्ग मापन-गुणक $\beta$
ज्ञात कीजिए, $y = \beta x + \varepsilon$ में। (ख)~अप्रणोदित अनुकरणों से आकलित उसकी मानक त्रुटि
\num{0.08} है। क्या अनुक्रिया संसूचित है? क्या वह मॉडल से संगत है? (ग)~उसी त्रुटि के साथ
0.5 के $\beta$ का क्या अर्थ होता?
\end{exercise}

\begin{exercise}[$\star\star\star$]\label{exo:b2:anthropogenic-climate-change:10}
दो अंगुलिछापें, ग्रीनहाउस गैसें $g = (0.2, 0.4, 0.6, 0.8)$ और एरोसॉल $a = (-0.1, -0.2,
-0.2, -0.3)$, दोनों चार दशकों में \unit{\celsius} में, प्रेक्षणों $y = (0.1,
0.2, 0.4, 0.5)$ पर समंजित की जाती हैं। (क)~न्यूनतम-वर्ग समंजन $y = \beta_g g + \beta_a a
+ \varepsilon$ के प्रसामान्य समीकरण लिखिए और उन्हें हल कीजिए। (ख)~अंतिम दशक में हर चालक पर आरोपित तापन ज्ञात कीजिए।
(ग)~जब दोनों अंगुलिछापें एक जैसी दिखती हैं तो $\beta_g$ और $\beta_a$ पर अनिश्चितताएँ क्यों बढ़ती हैं?
\end{exercise}

\begin{exercise}[$\star\star\star$]\label{exo:b2:anthropogenic-climate-change:11}
\omterm{def:b2:anthropogenic-climate-change:budget}{शेष कार्बन बजट} की विधि का उपयोग कीजिए, \qty{0.45}{\celsius} प्रति \qty{1000}{Gt} की TCRE,
2015 से 2019 तक उत्सर्जित \qty{210}{Gt} और पृथ्वी-तंत्र प्रतिपुष्टियों के लिए प्रति अंश \qty{26}{Gt} के साथ।
(क)~\qty{2}{\celsius} के लिए बजट ज्ञात कीजिए, जो 2010--2019 से \qty{0.93}{\celsius} ऊपर है,
यदि अन्य चालक \qty{0.1}{\celsius} जोड़ें। (ख)~अन्य चालकों का कितना अतिरिक्त तापन उसे आँके गए
\qty{1350}{Gt} के बराबर बना देता है? क्या वह आँके गए परास के भीतर है? (ग)~दिखाइए
कि हर \qty{0.1}{\celsius} गैर-CO$_2$ तापन किसी भी बजट को उतनी ही मात्रा से बदलता है।
\end{exercise}

\begin{exercise}[$\star\star\star$]\label{exo:b2:anthropogenic-climate-change:12}
शुद्ध प्रतिपुष्टि का बहुत संभावित परास $-1.81$ से \qty{-0.51}{W.m^{-2}.\celsius^{-1}} है, $-1.16$ के
आसपास। (क)~दोनों छोरों पर और केंद्र पर ECS ज्ञात कीजिए। (ख)~दिखाइए कि अपने केंद्र के परितः सममित $\lambda$
का परास ऐसा ECS परास देता है जो सममित नहीं है, और बताइए कि कौन-सी ओर लंबी है। (ग)~मूल्यांकनों को
कम ECS की तुलना में ऊँची ECS को ख़ारिज करना अधिक कठिन क्यों लगता है?
\end{exercise}
% ledger: ar6ch7:table710

\begin{problem}[{सप्ताहांत समस्या --- बजट को पढ़ना}]\label{pb:b2:anthropogenic-climate-change:1}
यह समस्या मूल्यांकन की ऊर्जा और कार्बन का प्रणोदन से \qty{1.5}{\celsius} के शेष बजट तक
अनुसरण करती है। अभ्यासों के आँकड़ों का उपयोग कीजिए, और: ERF 1750--2019 (\unit{W.m^{-2}})
कार्बन डाइऑक्साइड 2.16, मीथेन 0.54, नाइट्रस ऑक्साइड 0.21, हैलोजनित गैसें 0.41, ओज़ोन 0.47,
समतापमंडलीय जल-वाष्प 0.05, एरोसॉल--विकिरण $-0.22$, एरोसॉल--मेघ $-0.84$, भूमि-उपयोग
$-0.20$, हिम पर काला कार्बन 0.08, संघनन-पथ 0.06; कार्बन डाइऑक्साइड 2019 में \qty{410}{ppm} और
2025 में \qty{427.35}{ppm}; 2006--2018 पर तापन \qty{0.79}{W.m^{-2}}; ऊर्जा-बजट परिवर्तन
$\Delta N = 0.59$, $\Delta F = \qty{2.20}{W.m^{-2}}$, $\Delta T = \qty{1.03}{\celsius}$; TCRE
\qty{0.45}{\celsius} प्रति \qty{1000}{Gt} कार्बन डाइऑक्साइड (संभावित 0.27 से 0.63); 1850--2019 में उत्सर्जित
\qty{2390}{Gt}; 2010--2019 के लिए मानवीय तापन \qty{1.07}{\celsius}; 2015 से 2019 तक उत्सर्जित
\qty{210}{Gt}; अतिरिक्त गैर-CO$_2$ तापन \qty{0.1}{\celsius}, आँके गए 0.1
से \qty{0.2}{\celsius} का निचला छोर; पृथ्वी-तंत्र प्रतिपुष्टियाँ प्रति अंश \qty{26}{Gt}। तीन सार्थक
अंकों तक पूर्णांकित कीजिए।
% ledger: ipcc:erf-agents, ipcc:erf-2xco2, data:co2-mlo, ar6ch7:inventory, ar6ch7:energy-budget, ar6spm:tcre, ar6spm:budget, ar6ch5:budget-parts

\noindent\textbf{भाग I --- प्रणोदन।}
\begin{enumerate}
  \item चार सुमिश्रित ग्रीनहाउस गैसों के प्रणोदन का योग कीजिए।
  \item एरोसॉलों के प्रणोदन का योग कीजिए।
  \item सभी मानवीय पदों का योग कीजिए। आँके गए \qty{2.72}{W.m^{-2}} से तुलना कीजिए।
  \item लघुगणकीय नियम से 2019 में कार्बन डाइऑक्साइड के प्रणोदन की जाँच कीजिए।
  \item उसी नियम से 2025 में कार्बन डाइऑक्साइड का प्रणोदन ज्ञात कीजिए।
  \item किस पद की अनिश्चितता सबसे चौड़ी है, और वह कुल की अनिश्चितता पर क्यों हावी होती है?
\end{enumerate}

\noindent\textbf{भाग II --- ऊर्जा बजट।}
\begin{enumerate}[resume]
  \item 2006--2018 पर जलवायु तंत्र ने ज़ेटाजूल में कितनी ऊर्जा प्राप्त की?
  \item उसमें से कितनी महासागर में गई?
  \item ऊर्जा-बजट परिवर्तनों से प्रतिपुष्टि प्राचल ज्ञात कीजिए।
  \item प्रभावी \omterm{def:b2:anthropogenic-climate-change:ecs}{साम्य जलवायु संवेदनशीलता} ज्ञात कीजिए।
  \item उन्हीं आँकड़ों से \omterm{def:b2:anthropogenic-climate-change:tcr}{क्षणिक जलवायु अनुक्रिया} ज्ञात कीजिए।
  \item चित्र की शुद्ध प्रतिपुष्टि से ECS ज्ञात कीजिए, और दोनों मानों की तुलना कीजिए।
  \item भौतिक कारण बताइए कि ऐतिहासिक अनुमान निचला क्यों है।
\end{enumerate}

\noindent\textbf{भाग III --- प्रति टन तापन।}
\begin{enumerate}[resume]
  \item TCRE को कार्बन के प्रति \qty{1000}{PgC} अंशों में बदलिए।
  \item उसके संभावित परास को बदलिए।
  \item 1850--2019 के उत्सर्जनों को पेटाग्राम कार्बन में बदलिए।
  \item 1850--2019 के उत्सर्जनों के लिए TCRE कितना तापन देती है?
  \item \qty{1.07}{\celsius} के मानवीय तापन से तुलना कीजिए। यह तुलना अन्य मानवीय चालकों के
        शुद्ध प्रभाव के बारे में क्या बताती है?
\end{enumerate}

\noindent\textbf{भाग IV --- \qty{1.5}{\celsius} के लिए शेष बजट।}
\begin{enumerate}[resume]
  \item 2010--2019 से ऊपर कितना तापन शेष है?
  \item गैर-CO$_2$ तापन हटाइए और TCRE से भाग दीजिए: 2010--2019 के मध्य से कितना बजट
        शेष है?
  \item 2015--2019 के उत्सर्जन हटाइए।
  \item शेष तापन पर पृथ्वी-तंत्र प्रतिपुष्टियाँ हटाइए।
  \item यदि अन्य चालक इसके बजाय \qty{0.2}{\celsius} जोड़ें तो बजट कितना बदलता है?
  \item सारांश में: TCRE से परिकलित, 2020 के आरंभ से \qty{1.5}{\celsius} के लिए \omterm{def:b2:anthropogenic-climate-change:budget}{शेष कार्बन बजट}
        क्या है? आँके गए \qty{500}{Gt} से तुलना कीजिए।
\end{enumerate}
\end{problem}
